\chapter{Literature Review and Theoretical Foundations}
\label{ch:literature}

This chapter presents comprehensive review of existing research foundations underlying the two master problems formulated in the Introduction. The review synthesizes recent advances across seven interconnected research domains: adversarial machine learning for network security, graph neural networks for intrusion detection, federated learning under Byzantine adversaries, post-quantum cryptography and quantum-resistant security, state space models for temporal modeling, Bayesian deep learning and uncertainty quantification, and game-theoretic adversarial modeling. For each domain, we examine state-of-the-art approaches, identify fundamental limitations, and establish gaps motivating the novel contributions of this dissertation.

\section{Adversarial Machine Learning for Network Security}
\label{sec:lit_adversarial}

The application of deep learning to network intrusion detection has achieved remarkable performance on standard benchmarks, with recent models reporting accuracy exceeding 95\% on datasets including CIC-IDS2017, UNSW-NB15, and NSL-KDD. However, these models exhibit severe vulnerability to adversarial perturbations crafted through gradient-based optimization methods. Goodfellow et al. introduced the Fast Gradient Sign Method generating adversarial examples through single-step perturbations along the sign of the loss gradient. The method achieves high attack success rates with minimal distortion quantified through $\ell_p$ norms. Madry et al. proposed Projected Gradient Descent extending FGSM through iterative refinement with projection onto bounded perturbation sets defined by $\|\delta\|_\infty \leq \epsilon$, where $\delta$ denotes the perturbation vector and $\epsilon$ specifies the maximum allowed distortion.

Recent work has adapted these attacks to network intrusion detection contexts accounting for domain-specific constraints. Rigaki and Garcia demonstrated that adversarial examples generated for image classification translate to network traffic through feature space perturbations, achieving evasion rates of 73\% against deep neural network classifiers. Lin et al. introduced IDSGAN generating adversarial network traffic through generative adversarial networks trained to minimize detection probability while preserving attack functionality. The approach achieves 89\% evasion rate against state-of-the-art intrusion detection systems while maintaining successful exploitation of target vulnerabilities.

Defense mechanisms against adversarial attacks primarily focus on adversarial training incorporating perturbed examples during model training. Madry et al. showed that training on PGD-generated adversarial examples improves robust accuracy, defined as the fraction of correctly classified inputs under worst-case perturbations within radius $\epsilon$. However, adversarial training exhibits fundamental limitations characterized by the robustness-accuracy trade-off: improving robustness typically degrades standard accuracy on unperturbed test samples. Tsipras et al. provided theoretical analysis demonstrating that robust and accurate classifiers may not exist for certain data distributions, particularly when class boundaries require high-frequency components vulnerable to small perturbations.

Certified defense approaches provide formal robustness guarantees through conservative approximations of model behavior under perturbations. Cohen et al. introduced randomized smoothing achieving provable $\ell_2$ robustness through ensembling predictions on Gaussian-perturbed inputs. Given base classifier $f$, the smoothed classifier $g(x) = \argmax_c \Prob[f(x + \xi) = c]$ for $\xi \sim \mathcal{N}(0, \sigma^2 I)$ provides certified prediction constancy within radius $R$ computed from the probability gap between the top two classes. Wong and Kolter developed convex adversarial polytope analysis computing exact bounds on network output variation under perturbations, enabling certified robustness at computational cost scaling with network depth and width.

Despite these advances, current approaches exhibit critical limitations when applied to network intrusion detection. First, adversarial training achieves limited robustness gains while significantly increasing computational cost through iterative perturbation generation during training. Second, certified defense methods provide loose bounds resulting in low certified accuracy far below standard accuracy. Third, existing techniques focus primarily on $\ell_p$ threat models that may not align with realistic attack capabilities in network security contexts where adversaries face constraints on which features can be modified without invalidating attack payloads. Fourth, most work considers static threat models without accounting for adaptive adversaries that observe defense mechanisms and develop counter-strategies.

The gap between current adversarial robustness techniques and practical network security requirements motivates development of stochastic adversarial training methods providing stronger robustness through training on adversarial distributions rather than worst-case examples, combined with uncertainty quantification enabling rejection of predictions with high epistemic uncertainty.

\section{Graph Neural Networks for Intrusion Detection}
\label{sec:lit_gnn}

Graph neural networks have emerged as powerful tools for learning on graph-structured data through iterative message passing between connected nodes. Kipf and Welling introduced Graph Convolutional Networks aggregating neighbor features through spectral graph convolutions defined as $\mathbf{H}^{(l+1)} = \sigma(\tilde{\mathbf{D}}^{-1/2}\tilde{\mathbf{A}}\tilde{\mathbf{D}}^{-1/2}\mathbf{H}^{(l)}\mathbf{W}^{(l)})$, where $\tilde{\mathbf{A}} = \mathbf{A} + \mathbf{I}$ denotes the adjacency matrix with added self-connections, $\tilde{\mathbf{D}}$ is the degree matrix, $\mathbf{H}^{(l)}$ represents node features at layer $l$, and $\mathbf{W}^{(l)}$ contains learnable parameters. Veli\v{c}kovi\'{c} et al. proposed Graph Attention Networks incorporating learned attention mechanisms assigning importance weights to neighbors, enabling adaptive aggregation based on feature similarity.

Application of GNNs to network intrusion detection leverages graph representations of network traffic capturing relationships between hosts, services, and communication patterns. Lo et al. represented enterprise networks as attributed graphs where nodes correspond to hosts and edges represent network flows annotated with statistical features including packet counts, byte volumes, and inter-arrival times. GNN-based classification achieved 94.2\% accuracy detecting lateral movement attacks through learning propagation patterns indicative of reconnaissance and exploitation phases. Zhao et al. introduced heterogeneous GNNs operating on multi-relational graphs distinguishing different traffic types and protocol layers, achieving improved detection of application-layer attacks through explicit modeling of protocol dependencies.

Temporal graph neural networks extend static GNNs to time-evolving graphs through recurrent architectures or temporal attention mechanisms. Rossi et al. proposed Temporal Graph Networks maintaining dynamic node embeddings updated through temporal message passing as new edges arrive. The approach processes temporal events $\{(u_i, v_i, t_i, \mathbf{e}_i)\}$ representing interactions between nodes $u_i$ and $v_i$ at time $t_i$ with edge features $\mathbf{e}_i$, updating node memory states through gated recurrent units. Xu et al. introduced Inductive Representation Learning on Temporal Graphs combining temporal point processes with graph neural networks, modeling both continuous state evolution and discrete event occurrences through coupled differential equations.

Despite progress in GNN-based intrusion detection, current approaches face fundamental limitations restricting practical deployment. First, existing temporal GNNs rely on discrete time steps requiring choice of sampling interval. Too coarse sampling loses fine-grained temporal patterns while too fine sampling introduces computational overhead and redundant processing during inactive periods. Second, GNN architectures exhibit limited robustness to graph perturbations including edge additions, deletions, and feature modifications that adversaries can exploit. Third, scalability challenges arise when processing large-scale enterprise networks with millions of nodes and billions of edges, as message passing complexity scales quadratically with graph size. Fourth, most GNN approaches require labeled attack examples for supervised training, limiting ability to detect novel zero-day attacks absent from training data.

The limitations of discrete-time GNNs and their vulnerability to adversarial graph modifications motivate development of continuous-time temporal graph neural networks based on graph neural ordinary differential equations, enabling principled modeling of irregular temporal dynamics while providing robustness through smooth continuous state evolution. Additionally, the need for zero-shot generalization motivates integration of large language models with graph representations, enabling semantic reasoning about attack behaviors beyond pattern matching on labeled training examples.

\section{Federated Learning Under Byzantine Adversaries}
\label{sec:lit_federated}

Federated learning enables collaborative model training across distributed data sources without centralizing sensitive data, addressing privacy concerns in multi-organizational threat intelligence sharing. McMahan et al. introduced Federated Averaging aggregating local model updates from participating clients through weighted averaging proportional to dataset sizes. The protocol iterates between local training on client data for $E$ epochs using local learning rate $\eta_l$, and global aggregation computing $\mathbf{w}_{t+1} = \sum_{k=1}^K \frac{n_k}{n}\mathbf{w}_{t+1}^k$ where $\mathbf{w}_{t+1}^k$ denotes updated weights from client $k$ with $n_k$ samples and $n = \sum_k n_k$ total samples.

Privacy-preserving federated learning incorporates differential privacy guarantees limiting information leakage about individual training samples. Abadi et al. proposed Differentially Private Stochastic Gradient Descent clipping gradient norms to bound sensitivity and adding calibrated Gaussian noise $\mathcal{N}(0, \sigma^2 C^2)$ where $C$ denotes the clipping threshold and $\sigma$ is determined by privacy parameters $(\epsilon, \delta)$ through privacy amplification by sampling and composition theorems. Geyer et al. extended differential privacy to federated settings through local noise addition before aggregation, ensuring that aggregated model updates satisfy $(\epsilon, \delta)$-differential privacy even with untrusted aggregation servers.

Byzantine-robust federated learning addresses malicious participants submitting corrupted updates to degrade global model performance. Blanchard et al. introduced Krum aggregation selecting the update with smallest sum of squared distances to other updates, providing robustness when the fraction of Byzantine clients remains below $\frac{K-2}{K}$ where $K$ denotes total participants. Yin et al. proposed coordinate-wise trimmed mean and median aggregation removing extreme values across each parameter dimension, achieving convergence under $q < \frac{K}{2}$ Byzantine clients with convergence rate $O\left(\sqrt{\frac{q}{K}}\right)$.

Optimal transport theory provides principled frameworks for aligning distributions across federated clients exhibiting heterogeneous data distributions. Redko et al. applied optimal transport to domain adaptation computing transport plans $\mathbf{T}^* \in \argmin_{\mathbf{T} \in \Pi(\mu_s, \mu_t)} \langle \mathbf{C}, \mathbf{T} \rangle$ minimizing expected transport cost between source distribution $\mu_s$ and target distribution $\mu_t$ subject to marginal constraints $\Pi(\mu_s, \mu_t)$, where cost matrix $\mathbf{C}$ quantifies dissimilarity between samples. Courty et al. introduced Joint Distribution Optimal Transport learning feature representations minimizing both classification loss and optimal transport distance between domains.

Despite advances in Byzantine-robust and privacy-preserving federated learning, significant gaps remain for network security applications. First, existing Byzantine-robust aggregation rules require knowing or estimating the fraction of malicious clients, which may be difficult in practice as adversaries can adaptively adjust attack strategies. Second, the privacy-utility trade-off under differential privacy remains unfavorable, with strong privacy guarantees requiring noise levels that degrade detection accuracy below operationally acceptable thresholds. Third, optimal transport-based domain adaptation has not been integrated with differential privacy, limiting applicability to privacy-sensitive security contexts. Fourth, federated intrusion detection systems lack integration with large language models that could enable semantic understanding of attack patterns and zero-shot generalization to novel threats.

These limitations motivate development of integrated federated learning frameworks combining Byzantine-robust aggregation, differentially private optimal transport for distribution alignment, and large language model integration for zero-shot detection capabilities, while providing formal convergence guarantees under adversarial conditions.

\section{Post-Quantum Cryptography and Quantum-Resistant Security}
\label{sec:lit_postquantum}

The advancement of quantum computing threatens security foundations of current cryptographic protocols through Shor's algorithm enabling efficient factorization and discrete logarithm computation. The National Institute of Standards and Technology initiated standardization of post-quantum cryptographic algorithms resistant to quantum attacks, with standardized algorithms announced in 2024 including CRYSTALS-Kyber for key encapsulation, CRYSTALS-Dilithium for digital signatures, and Sphincs+ for hash-based signatures. These algorithms provide security foundations for quantum era but introduce new attack surfaces requiring adapted intrusion detection capabilities.

Lattice-based cryptography including CRYSTALS-Kyber and Dilithium provides security guarantees based on hardness of Learning With Errors problem. Given matrix $\mathbf{A} \in \mathbb{Z}_q^{m \times n}$ and target vector $\mathbf{b} = \mathbf{A}\mathbf{s} + \mathbf{e} \bmod q$ for secret $\mathbf{s}$ and small error $\mathbf{e}$, the LWE problem requires recovering $\mathbf{s}$ which remains hard even for quantum adversaries when parameters satisfy appropriate security bounds. However, implementation vulnerabilities including timing side-channels and fault injection attacks exploit correlations between execution time or power consumption and secret values. Ravi et al. demonstrated timing attacks against Kyber implementations recovering secret keys through statistical analysis of decapsulation timing variations correlated with Hamming weights of intermediate values.

Code-based cryptography provides alternative post-quantum foundation through hardness of decoding random linear codes. McEliece and Niederreiter cryptosystems achieve security from computational intractability of syndrome decoding problem. Despite strong theoretical security, practical implementations face challenges including large public key sizes exceeding 100 KB and vulnerability to side-channel attacks exploiting patterns in decoding algorithms. Structural attacks against specific parameter choices have been demonstrated, requiring careful parameter selection validated through extensive cryptanalysis.

Hash-based signatures including Sphincs+ provide conservative security foundations based solely on collision-resistant hash functions without additional hardness assumptions. The schemes achieve stateless signature generation through randomization avoiding key exhaustion vulnerabilities of earlier hash-based approaches. However, large signature sizes and computational overhead during signing operation pose practical deployment challenges requiring efficiency optimizations.

Current intrusion detection systems lack capabilities to detect attacks specifically targeting post-quantum cryptographic implementations. Existing anomaly detection approaches trained on classical cryptographic protocols may fail to generalize to post-quantum traffic patterns exhibiting different timing distributions, message sizes, and protocol flows. Side-channel attacks against post-quantum implementations require specialized detection mechanisms analyzing timing patterns and power consumption characteristics distinct from classical cryptography. Additionally, the transition period where systems support both classical and post-quantum algorithms introduces complexity where adversaries may force downgrade attacks to weaker classical protocols.

These gaps motivate development of specialized post-quantum intrusion detection systems incorporating dedicated feature extraction for lattice-based, code-based, and hash-based cryptographic protocols, with detection capabilities for timing side-channels, structural attacks, and downgrade attempts. Integration with continuous-time temporal modeling enables detection of subtle timing variations indicative of side-channel exploitation without requiring fixed sampling intervals.

\section{State Space Models for Temporal Sequence Modeling}
\label{sec:lit_ssm}

State space models provide elegant mathematical frameworks for modeling temporal dynamics through linear differential equations. The continuous-time state space representation $\frac{d\mathbf{x}(t)}{dt} = \mathbf{A}\mathbf{x}(t) + \mathbf{B}\mathbf{u}(t)$ and observation equation $\mathbf{y}(t) = \mathbf{C}\mathbf{x}(t) + \mathbf{D}\mathbf{u}(t)$ capture evolution of latent state $\mathbf{x}(t) \in \mathbb{R}^N$ driven by input $\mathbf{u}(t)$ and mapped to output $\mathbf{y}(t)$ through projection matrix $\mathbf{C}$. Classical Kalman filtering provides optimal state estimation under Gaussian assumptions, while variants including Extended Kalman Filter and Unscented Kalman Filter handle nonlinear dynamics.

Recent deep learning approaches combine state space models with neural networks enabling learning of complex nonlinear dynamics. Gu et al. introduced Structured State Space Sequence models achieving efficient sequence modeling through diagonal plus low-rank parameterization of state transition matrix $\mathbf{A}$ enabling fast computation through Fast Fourier Transform. The approach achieves performance competitive with Transformers on long-range sequence tasks while maintaining linear computational complexity in sequence length compared to quadratic complexity of attention mechanisms.

Mamba, introduced by Gu and Dao in 2024, extends structured state spaces through selective scan mechanism making state transition matrices input-dependent. Unlike fixed SSMs with time-invariant $\mathbf{A}$, Mamba computes $\mathbf{A}(t) = f_A(\mathbf{u}(t))$ and $\mathbf{B}(t) = f_B(\mathbf{u}(t))$ through learned projections of input features. This selection mechanism enables adaptive filtering of relevant information while maintaining efficient parallel computation through hardware-aware algorithms. Mamba achieves state-of-the-art performance on language modeling benchmarks including Pile and C4 datasets, outperforming Transformers of comparable size while exhibiting linear complexity in sequence length.

Application of state space models to network intrusion detection remains limited despite their theoretical advantages for temporal modeling. Lo et al. applied classical Kalman filtering to track network flow statistics, detecting anomalies through innovation sequences measuring discrepancy between predicted and observed measurements. However, linear Kalman filtering cannot capture complex nonlinear attack patterns. Neural state space models including Mamba have not been systematically explored for security applications despite potential advantages including efficient long-range temporal modeling, selective attention to relevant features, and parameter efficiency compared to Transformer architectures.

Furthermore, robustness of state space models under adversarial perturbations and data poisoning attacks has received limited attention. Training on corrupted data containing mislabeled samples or manipulated features can degrade model performance, but mechanisms for ensuring robust learning under poisoning remain underdeveloped. PAC-Bayesian analysis provides theoretical frameworks for generalization under distribution shift but has not been integrated with state space models to provide formal robustness certificates.

These limitations motivate development of adversarially resilient state space models specifically designed for intrusion detection, incorporating progressive adversarial robustness distillation to ensure learning from potentially poisoned data, and PAC-Bayesian analysis providing formal generalization bounds supporting deployment in safety-critical security contexts.

\section{Bayesian Deep Learning and Uncertainty Quantification}
\label{sec:lit_bayesian}

Bayesian deep learning provides principled frameworks for uncertainty quantification through probabilistic inference over model parameters. Given training data $\mathcal{D} = \{(\mathbf{x}_i, y_i)\}_{i=1}^n$, Bayesian inference computes posterior distribution $p(\theta|\mathcal{D}) = \frac{p(\mathcal{D}|\theta)p(\theta)}{\int p(\mathcal{D}|\theta)p(\theta)d\theta}$ over parameters $\theta$ combining likelihood $p(\mathcal{D}|\theta)$ and prior $p(\theta)$. Prediction incorporates parameter uncertainty through marginalization: $p(y^*|\mathbf{x}^*, \mathcal{D}) = \int p(y^*|\mathbf{x}^*, \theta)p(\theta|\mathcal{D})d\theta$.

Practical Bayesian inference employs approximate methods due to intractability of exact posterior computation for neural networks. Variational inference approximates posterior through optimization over parametric distribution family $q_\phi(\theta)$ minimizing KL divergence $KL(q_\phi(\theta)||p(\theta|\mathcal{D})) = \mathbb{E}_{q_\phi}[\log q_\phi(\theta) - \log p(\theta|\mathcal{D})]$, equivalently maximizing evidence lower bound $\mathcal{L}_{ELBO} = \mathbb{E}_{q_\phi}[\log p(\mathcal{D}|\theta)] - KL(q_\phi(\theta)||p(\theta))$. Blundell et al. introduced Bayes by Backprop implementing variational inference through reparameterization trick enabling gradient-based optimization of variational parameters $\phi$.

Monte Carlo Dropout provides computationally efficient uncertainty estimation through stochastic regularization. Gal and Ghahramani showed that dropout training approximates variational inference over model parameters, enabling uncertainty quantification by performing multiple forward passes with different dropout masks and computing prediction variance. Given model $f_w$ with dropout, epistemic uncertainty is estimated as $\mathbb{V}[\hat{y}] = \frac{1}{T}\sum_{t=1}^T f_{w_t}(\mathbf{x})^2 - (\frac{1}{T}\sum_{t=1}^T f_{w_t}(\mathbf{x}))^2$ where $w_t$ denotes sampled weight configurations.

Uncertainty decomposition distinguishes epistemic uncertainty arising from model uncertainty reducible through additional training data, and aleatoric uncertainty representing inherent data noise irreducible through more data. Kendall and Gal showed that aleatoric uncertainty can be modeled through additional output heads predicting observation noise, enabling learned heteroscedastic noise estimation. Total predictive uncertainty decomposes as $\mathbb{U}_{total} = \mathbb{U}_{epistemic} + \mathbb{U}_{aleatoric}$.

Calibration measures alignment between predicted confidence and empirical accuracy. Guo et al. demonstrated that modern neural networks produce poorly calibrated predictions with confidence exceeding accuracy. Expected Calibration Error quantifies miscalibration by binning predictions by confidence and measuring average absolute difference between confidence and accuracy: $ECE = \sum_{m=1}^M \frac{|B_m|}{n}|\text{acc}(B_m) - \text{conf}(B_m)|$ where $B_m$ denotes the $m$-th confidence bin. Temperature scaling provides simple post-hoc calibration through learned temperature parameter $T$ in softmax: $\hat{p}_i = \frac{\exp(z_i/T)}{\sum_j \exp(z_j/T)}$ optimized to minimize validation set ECE.

Despite these advances, integration of Bayesian uncertainty quantification with adversarial robustness remains limited. Existing work focuses primarily on either uncertainty estimation or adversarial training, but unified frameworks addressing both objectives simultaneously are rare. Additionally, uncertainty-aware adversarial training could enable adaptive attack strategies where adversaries target samples with high epistemic uncertainty exploiting model lack of knowledge. Theoretical characterization of uncertainty under adversarial perturbations requires formal analysis of how input modifications propagate through probabilistic neural networks to affect posterior predictive distributions.

Application of Bayesian methods to network intrusion detection remains limited despite potential benefits including quantified confidence for security-critical decisions, principled handling of concept drift through online Bayesian updating, and uncertainty-guided active learning for efficient labeling of ambiguous samples. Current intrusion detection systems provide point predictions without confidence estimates, limiting ability to appropriately triage alerts and allocate analyst investigation resources.

These gaps motivate development of Bayesian deep learning frameworks specifically designed for adversarially robust intrusion detection, incorporating variational attention mechanisms enabling uncertainty quantification while maintaining computational tractability, stochastic adversarial training providing robustness under attack, and proper calibration ensuring that predicted confidences align with empirical accuracy enabling reliable security decision making.

\section{Game-Theoretic Adversarial Modeling}
\label{sec:lit_game_theory}

Game-theoretic frameworks model strategic interactions between defenders deploying security mechanisms and attackers optimizing exploit strategies. Stackelberg game formulations capture asymmetric information and sequential decision making where defenders commit to strategies observable by attackers who then optimize responses. The defender solves $\max_{\pi_d} \min_{\pi_a} U_d(\pi_d, \pi_a)$ where $\pi_d$ denotes defender strategy (e.g., deployed detection model), $\pi_a$ represents attacker strategy (e.g., evasion technique), and $U_d$ specifies defender utility typically combining detection accuracy and operational costs.

Nash equilibrium characterizes stable strategy profiles where no player benefits from unilateral deviation. In symmetric zero-sum security games, maximin and minimax strategies coincide at the Nash equilibrium value. However, realistic security scenarios often involve asymmetric payoffs, incomplete information about adversary capabilities, and sequential rather than simultaneous moves requiring extended equilibrium concepts including subgame perfect equilibrium and Bayesian Nash equilibrium incorporating player beliefs.

Adversarial classification games model attackers manipulating features to evade detection while defenders design robust classifiers accounting for strategic manipulation. Brückner and Scheffer formulated classifier learning as game between learner and data generator, deriving Nash equilibrium classifiers robust to anticipated feature manipulation. Globerson and Roweis considered worst-case feature deletion attacks optimizing linear classifiers maximizing margin under worst-case test-time deletions.

Online learning frameworks enable adaptive strategy selection without requiring prior knowledge of opponent strategies. Hannan-consistent algorithms including multiplicative weights update maintain sublinear regret $R(T) = o(T)$ meaning average per-round regret vanishes: $\frac{R(T)}{T} \to 0$. No-regret learning algorithms guarantee convergence to Nash equilibrium in repeated zero-sum games. Freund and Schapire introduced Hedge algorithm maintaining weight distribution over actions with multiplicative updates proportional to exponential of cumulative payoffs, achieving regret bound $O(\sqrt{T \log K})$ for $K$ actions over $T$ rounds.

Recent work integrates game theory with machine learning for adaptive adversarial defense. Liu et al. formulated adversarial example generation as attacker strategy in zero-sum game, deriving defense mechanisms through minimax optimization. Madry et al.'s PGD adversarial training corresponds to empirical risk minimization under worst-case perturbations, interpretable as approximate solution to defender's maximin problem. However, fixed perturbation budgets fail to capture realistic attacker cost structures where different feature modifications incur heterogeneous costs based on domain constraints.

Limitations of current game-theoretic approaches for network security include simplified threat models not capturing full range of attacker capabilities and constraints, computational intractability of computing exact equilibria for large strategy spaces, and lack of integration with uncertainty quantification enabling risk-aware decision making under incomplete information. Most game-theoretic formulations assume common knowledge of payoff structures and strategy spaces which may not hold in security contexts with novel attacks and unknown adversary objectives.

Furthermore, online learning algorithms providing no-regret guarantees require carefully designed loss feedback. In security applications, true attack outcomes may be delayed, censored, or manipulated by adversaries, violating standard online learning assumptions. Bandit feedback scenarios where defenders observe only outcomes of deployed strategies rather than full loss information require modified algorithms with weaker regret bounds.

These limitations motivate development of game-theoretic frameworks specifically designed for adversarial network security incorporating realistic cost models reflecting domain constraints on feature manipulation, integration with Bayesian uncertainty quantification enabling principled decision making under incomplete information, and online learning algorithms providing provable regret bounds under adversarial feedback where outcomes may be delayed or manipulated by strategic attackers.

\section{Identified Research Gaps}
\label{sec:lit_gaps}

The comprehensive literature review reveals critical gaps between current state of the art and practical requirements for adversarially resilient intrusion detection in heterogeneous network environments:

\textbf{Gap 1: Insufficient Adversarial Robustness with Formal Guarantees.} Existing adversarial training methods achieve limited robustness improvements at significant computational cost without providing formal certificates on worst-case performance. Certified defense approaches provide loose bounds resulting in low certified accuracy substantially below standard accuracy. Network security applications require both high empirical robustness against observed attacks and formal guarantees supporting safety-critical deployment in adversarial environments.

\textbf{Gap 2: Unfavorable Privacy-Utility Trade-offs in Federated Learning.} Differentially private federated learning exhibits severe accuracy degradation under privacy parameters meeting regulatory requirements. Byzantine-robust aggregation rules require knowing or estimating the fraction of malicious participants which may be difficult when adversaries adapt strategies dynamically. Integration of differential privacy with optimal transport-based domain adaptation remains unexplored despite potential for improving cross-domain generalization under privacy constraints.

\textbf{Gap 3: Limited Cross-Domain Generalization Capabilities.} Models trained on specific network environments fail to generalize to deployment contexts exhibiting different traffic characteristics, topology structures, and attack manifestations. Existing domain adaptation techniques require target domain data for fine-tuning, limiting zero-shot generalization capabilities needed for proactive defense against novel attack vectors. Large language models provide semantic reasoning capabilities that could enable understanding of attack principles, but integration with temporal graph representations of network traffic remains unexplored.

\textbf{Gap 4: Discrete-Time Limitations of Temporal Models.} Current temporal graph neural networks rely on discrete time steps requiring sampling interval selection that trades off between temporal resolution and computational efficiency. Irregular network events occurring at multiple time scales from microseconds to months cannot be efficiently modeled through fixed sampling. Continuous-time modeling through neural ordinary differential equations provides mathematical elegance but has not been systematically applied to graph-structured network security data.

\textbf{Gap 5: Absence of Post-Quantum Attack Detection Capabilities.} Migration to post-quantum cryptographic protocols introduces new attack surfaces including timing side-channels against lattice-based schemes and structural attacks on code-based cryptography. Existing intrusion detection systems trained on classical cryptographic traffic patterns lack specialized feature extraction and detection capabilities for post-quantum protocols. The transition period supporting both classical and post-quantum algorithms enables downgrade attacks requiring detection mechanisms identifying forced algorithm rollback attempts.

\textbf{Gap 6: Limited Integration of Uncertainty Quantification with Robustness.} Bayesian deep learning provides principled uncertainty quantification but lacks integration with adversarial robustness objectives. Uncertainty-aware adversarial training could enable adaptive defense strategies rejecting predictions with high epistemic uncertainty for human analyst review, but theoretical frameworks characterizing uncertainty under adversarial perturbations remain underdeveloped. Proper calibration ensuring alignment between predicted confidence and empirical accuracy is essential for security-critical decision making but receives limited attention in adversarial robustness literature.

\textbf{Gap 7: Simplified Threat Models in Game-Theoretic Frameworks.} Current game-theoretic adversarial modeling employs simplified assumptions including fixed perturbation budgets not reflecting realistic cost structures, complete information about opponent capabilities contradicting incomplete knowledge in security contexts, and failure to integrate with online learning algorithms providing adaptive strategy selection under non-stationary adversarial behavior.

These identified gaps directly motivate the master problem formulations presented in the Introduction and establish requirements for the novel theoretical frameworks, algorithmic implementations, and empirical validations developed in subsequent chapters. The dissertation addresses these gaps through integrated solutions combining continuous-time graph neural networks, optimal transport-based federated learning, post-quantum protocol analysis, state space models with adversarial robustness, Bayesian uncertainty quantification, and game-theoretic online learning providing comprehensive treatment of adversarial resilience in hybrid intrusion detection systems.
